\subsection{两个集合的并集与交集}

若 A、B 是集合，我们记

\[
\mathrm{A} \cup \mathrm{B} = \bigcup_ {\mathbf {X} \in \{\mathbf {A}, \mathrm{B} \}} \mathrm{X}, \quad \mathrm{A} \cap \mathrm{B} = \bigcap_ {\mathbf {X} \in \{\mathbf {A}, \mathrm{B} \}} \mathrm{X}.
\]

显然，A ∪ B 是所有属于 A 或属于 B（或同时属于两者）的对象的集合，而 A ∩ B 是所有同时属于 A 和 B 的对象的集合。特别地， \{x, y\} = \{x\} ∪ \{y\}. 设 \{x, y, z\} = \{x, y\} ∪ \{z\}。集合 \{x, y, z\} 是仅以 x、y 和 z 为元素的集合。类似地，我们写

\[
\{x, y, z, t \} = \{x, y, z \} \cup \{t \},
\]

等等。

如果 A、B、C、D 是集合，我们写

\[
\begin{array}{c} \mathrm {A \cup B \cup C = \bigcup_ {\mathbf {X} \in \{\mathrm{A,B,C} \}} X, \quad A \cap B \cap C = \bigcap_ {\mathbf {X} \in \{\mathrm{A,B,C} \}} X;} \\ \mathrm {A \cup B \cup C \cup D = \bigcup_ {\mathbf {X} \in \{\mathrm{A,B,C,D} \}} X, \quad A \cap B \cap C \cap D = \bigcap_ {\mathbf {X} \in \{\mathrm{A,B,C,D} \}} X;} \end{array}
\]

等等。

设 A、B、C 是集合。由命题 1 和命题 2 我们推出公式

\[
\begin{array}{c} \mathbf {A} \cup \mathbf {B} = \mathbf {B} \cup \mathbf {A}, \qquad \mathbf {A} \cap \mathbf {B} = \mathbf {B} \cap \mathbf {A}, \\ \mathbf {A} \cup (\mathbf {B} \cup \mathbf {C}) = (\mathbf {A} \cup \mathbf {B}) \cup \mathbf {C} = \mathbf {A} \cup \mathbf {B} \cup \mathbf {C}, \\ \mathbf {A} \cap (\mathbf {B} \cap \mathbf {C}) = (\mathbf {A} \cap \mathbf {B}) \cap \mathbf {C} = \mathbf {A} \cap \mathbf {B} \cap \mathbf {C}. \end{array}
\]

这些公式也是准则 C24（第一章，§3，第 5 号）中所陈述定理的直接推论。类似地，可以证明公式

\[
\mathbf {A} \cup (\mathbf {B} \cap \mathbf {C}) = (\mathbf {A} \cup \mathbf {B}) \cap (\mathbf {A} \cup \mathbf {C}), \quad \mathbf {A} \cap (\mathbf {B} \cup \mathbf {C}) = (\mathbf {A} \cap \mathbf {B}) \cup (\mathbf {A} \cap \mathbf {C})
\]

（并关于交的“分配律”以及交关于并的分配律；参见 §5，第 6 号）。

关系 \(\mathbf{A} \subset \mathbf{B}\) 等价于 \(\mathbf{A} \cup \mathbf{B} = \mathbf{B}\) 以及 \(\mathbf{A} \cap \mathbf{B} = \mathbf{A}\) 。如果 \(\mathbf{A}\) 和 \(\mathbf{B}\) 是集合 \(\mathbf{E}\) 的子集，我们从命题5（或从准则C24）推导出公式

\[
\mathbb {G} _ {\mathbf {E}} (\mathbf {A} \cup \mathbf {B}) = \left(\mathbb {G} _ {\mathbf {E}} \mathbf {A}\right) \cap \left(\mathbb {G} _ {\mathbf {E}} \mathbf {B}\right), \quad \mathbb {G} _ {\mathbf {E}} (\mathbf {A} \cap \mathbf {B}) = \left(\mathbb {G} _ {\mathbf {E}} \mathbf {A}\right) \cup \left(\mathbb {G} _ {\mathbf {E}} \mathbf {B}\right);
\]

此外，我们有

\[
\mathrm{A} \cup (\mathbb {C} _ {\mathrm{E}} \mathrm{A}) = \mathrm{E}, \quad \mathrm{A} \cap (\mathbb {C} _ {\mathrm{E}} \mathrm{A}) = \varnothing .
\]

如果 \(\Gamma\) 是...之间的对应关系 \(E\) 并且 \(F\) ，并且如果 \(A, B\) 是……的子集 \(E\) ，由命题3可知

\[
\Gamma \langle \mathrm{A} \cup \mathrm{B} \rangle = \Gamma \langle \mathrm{A} \rangle \cup \Gamma \langle \mathrm{B} \rangle , \quad \Gamma \langle \mathrm{A} \cap \mathrm{B} \rangle \subset \Gamma \langle \mathrm{A} \rangle \cap \Gamma \langle \mathrm{B} \rangle ,
\]

并且，如果 \(f\) 是从 \(F\) 到 \(E\) 的映射，

\[
\overline {{{f}}} ^ {1} \langle \mathrm{A} \cap \mathrm{B} \rangle = \overline {{{f}}} ^ {1} \langle \mathrm{A} \rangle \cap \overline {{{f}}} ^ {1} \langle \mathrm{B} \rangle
\]

由命题4得出。

我们还记录以下关于补集的命题：

命题6. 设 \(f\) 是从 \(A\) 到 \(B\) 的映射。对于 \(B\) 的每一个子集 \(Y\) ，都有

\[
\overline {{{f}}} ^ {1} \langle \mathrm{B} - \mathrm{Y} \rangle = \overline {{{f}}} ^ {1} \langle \mathrm{B} \rangle - \overline {{{f}}} ^ {1} \langle \mathrm{Y} \rangle .
\]

因为 \(x\) 属于 \(\overline{f}^1\langle \mathbf{B} - \mathbf{Y}\rangle\) 当且仅当 \(f(x)\) 属于 \(\mathbf{B}\) 而不属于 \(\mathbf{Y}\) ，也就是当且仅当 \(x\) 属于 \(\overline{f}^1\langle \mathbf{B}\rangle\) 而不属于 \(\overline{f}^1\langle \mathbf{Y}\rangle\) .

推论。设 \(f\) 是从 \(\mathbf{A}\) 到 \(\mathbf{B}\) 的单射。对于 \(\mathbf{A}\) 的每一个子集 \(\mathbf{X}\) ，都有 \(f\langle \mathbf{A} - \mathbf{X}\rangle = f\langle \mathbf{A}\rangle - f\langle \mathbf{X}\rangle\) .

写成 \(f = i \circ g\) ，其中 \(i\) 是从 \(f\langle A\rangle\) 到 B 的典范单射，就把该推论化归为对 \(\overline{g}^1\) 应用命题6。

¶ 交集 X ∩ A 有时称为 X 在 A 上的迹。若 F 是一个集合族，则 F 中集合在 A 上的迹所组成的集合称为 F 在 A 上的迹。

